% Sumber persamaan Bab III; nomor akhir mengikuti naskah utama.
% Persamaan 3.1
\[
C_e=0{,}55H_e+0{,}30S_e+0{,}15W_e
\]

% Persamaan 3.2
\[
C_e=0{,}55(1)+0{,}30(1)+0{,}15(1)=1{,}00
\]

% Persamaan 3.3
\[
C_e=0{,}55(-10)+0{,}30(-0{,}5)+0{,}15(-0{,}5)=-5{,}725
\]

% Persamaan 3.4
\[
b=\min\left(24,\max\left(0,\left\lfloor\frac{d}{6000}\times25\right\rfloor\right)\right)
\]

% Persamaan 3.5
\[
D_t=d_H(v_t,v_0)
\]

% Persamaan 3.6
\[
R_e=C_e\frac{L_e}{100}
\]

% Persamaan 3.7
\[
R_a=-0{,}10\frac{L_e}{100}I_p
\]

% Persamaan 3.8
\[
R_d=1{,}20\frac{D_{\mathrm{sebelum}}-D_{\mathrm{sesudah}}}{100}
\]

% Persamaan 3.9
\[
R_\delta=-5\frac{|d_{\mathrm{akhir}}-5000|}{100}
\]

% Persamaan 3.10
\[
r=R_e+R_a+R_d+R_v+R_T+R_\delta
\]

% Persamaan 3.11
\[
Q(s,a)\leftarrow Q(s,a)+\alpha\left[r+\gamma\max_{a^\prime\in A_{\mathrm{boleh}}}Q(s^\prime,a^\prime)-Q(s,a)\right]
\]

% Persamaan 3.12
\[
\varepsilon_n=0{,}05+0{,}95\left[1-\min\left(\frac{n-1}{D-1},1\right)\right]
\]

% Persamaan 3.13
\[
\varepsilon_{50.000}=0{,}05+0{,}95\left(1-\frac{49.999}{99.999}\right)\approx0{,}5250
\]

% Persamaan 3.14
\[
I_{\mathrm{loop}}=\mathbf{1}\{v_{\mathrm{akhir}}=v_0\;\land\;4000\le d_{\mathrm{akhir}}\le6000\}
\]

% Persamaan 3.15
\[
O=\frac{N-U}{N}
\]

% Persamaan 3.16
\[
\overline C=\frac{\sum_{i=1}^{N}C_iL_i}{\sum_{i=1}^{N}L_i}
\]

% Persamaan 3.17
\[
P=\frac{n_{\mathrm{mendekat}}}{n_{\mathrm{pulang}}}
\]

% Persamaan 3.18
\[
S=\frac{n_{\mathrm{berhasil}}}{n_{\mathrm{dinilai}}}\times100\%
\]

% Persamaan 3.19
\[
T_j=\frac{n_j}{n_{\mathrm{dinilai}}}\times100\%
\]

% Persamaan 3.20
\[
d_{\mathrm{akhir}}=1200+2000+1650=4850\;\mathrm{m}
\]

% Persamaan 3.21
\[
E=|4850-5000|=150\;\mathrm{m}
\]

% Persamaan 3.22
\[
R_{\mathrm{total}}=-12{,}50+100-10=77{,}50
\]

% Persamaan 3.23
\[
\overline C=\frac{1000(1)+3000(0{,}325)+1000(0{,}725)}{5000}=0{,}54
\]

% Persamaan 3.24
\[
O=\frac{10-8}{10}=0{,}20=20\%
\]
